\chapter{Contexte general du projet}

\section{Introduction}
Ce chapitre presente le cadre general de notre projet, l'organisme d'accueil, l'analyse de l'existant, la problematique et la solution proposee.

\section{Presentation de l'organisme d'accueil}
\textit{[A completer avec les informations exactes de l'entreprise d'accueil : historique, missions, organisation, departement d'affectation.]}

\section{Contexte du domaine}
Le secteur hotelier tunisien connait une evolution importante en matiere de digitalisation. Les etablissements cherchent a ameliorer leur visibilite, leur qualite de service et leur capacite a repondre rapidement aux attentes des touristes nationaux et internationaux.

\section{Analyse de l'existant}
\subsection{Etude de l'existant}
Les approches traditionnelles reposent souvent sur des FAQ statiques, des supports disperses et une interaction limitee avec les utilisateurs.

\subsection{Problematique}
Les principales limites identifiees sont :
\begin{itemize}
  \item absence d'assistance conversationnelle contextuelle pour chaque hotel ;
  \item manque de personnalisation selon la langue et le profil client ;
  \item faible exploitation analytique des interactions utilisateurs ;
  \item difficultes de gestion centralisee des contenus hoteliers.
\end{itemize}

\subsection{Solution proposee}
Nous proposons une plateforme web intelligente qui integre un chatbot RAG multilingue, un tableau de bord analytics et un espace administrateur securise pour gerer les donnees hotelieres.

\section{Methodologie et planification}
\textit{[A completer : methode de gestion de projet choisie (Scrum/Kanban), planning, jalons, livrables.]}

\section{Conclusion}
Ce chapitre a pose le contexte du projet et justifie la solution retenue. Le chapitre suivant detaille les besoins fonctionnels et non fonctionnels ainsi que l'architecture cible.
